\section{Discussion and limitations}
\label{sec:discussion}

\subsection{What the design buys, and what it costs}
\label{sec:tradeoffs}

The three mechanisms are separable in principle but reinforce each other in
practice.  Open oracles without structure preservation would be a typing
game on top of circuits that still explode with scale; structure
preservation without protocols would preserve modules no library can
assemble; protocols without provenance would check shapes while losing the
mathematics.  The costs are equally real.  Embedding in Python adopts the
host's dynamic typing, so contract checks run at generation time rather
than at a static pass; the IR's concreteness (all widths and angles
resolved) rules out symbolic parameters inside RIR, pushing parameter
studies to the generation layer; and the unitary-only IR, while it keeps
solver outputs composable, means hybrid quantum--classical algorithms
orchestrate from the host rather than from the IR---a deliberate division
that favors the fault-tolerant algorithm regime over the variational one.

\subsection{Limitations of the current implementation}
\label{sec:limitations}

We enumerate the limitations most relevant to prospective users and to the
future-validation program of \secref{sec:reserved}:

\begin{itemize}
\item \textbf{Solver certification is small-scale.}  All six ODE families,
  both QLSS kernels, QHAM, and QFVM are certified against truncated-oracle
  references at register-exact scale (\secref{sec:numverify}), and each
  still carries recorded pending items (finite quadrature, omitted contour
  poles, Carleman tails, recovery windows, success channels); the
  Schr\"odingerization momentum-sign and oblivious-amplification iteration
  defects found by the campaign are repaired and pinned by regression
  tests.  None is
  a certified production solver yet (Table~\ref{tab:status-matrix}).
\item \textbf{Hamiltonian-simulation kernels are partial.}  First-order
  Trotter and truncated Taylor are built in; a general QSP-based
  Hamiltonian-simulation kernel must currently be \emph{injected} by the
  caller (the QSVT-based block-encoding route exists separately).
  Higher-order product formulas are not yet provided.
\item \textbf{Amplitude-estimation variants are canonical only.}  No
  maximum-likelihood or robust amplitude estimation; no VTAA variants for
  the QLSS kernels.
\item \textbf{Width regime.}  Single registers span at most 64 bits (a
  backend storage constraint); wider objects are fused multi-register
  views.  Explicit QHAM tensor lifts are capped by this width.
\item \textbf{Number theory is illustrative.}  The modular-multiplication
  path is a bounded permutation synthesis for small registers, not scalable
  Shor arithmetic---and is labeled as such.
\item \textbf{No physical-layer resource accounting}---the logical
  Toffoli+Clifford+T+QRAM estimator of \secref{sec:resources} does not yet
  model synthesis passes, depth, distillation factories, or surface-code
  cycles; scoped in \secref{sec:reserved}.
\end{itemize}

\subsection{Future work}
\label{sec:future}

Beyond the reserved validation axis (backend campaigns at scale, solver
certification studies, and a first-class resource estimator), the natural
extensions are: a general QSP Hamiltonian-simulation kernel closing the
injection point; vendoring a proven high-degree phase-synthesis backend
(e.g.\ Laurent exact decomposition) behind the existing phase-value
interface---an implementation upgrade that, by design, requires no change
to the language, the IR, or the verification pass; robust and
maximum-likelihood amplitude estimation; a
wider nonlinear-PDE repertoire on the QHAM path (systems, higher
dimensions); classical-optimization protocol objects for variational
workloads; and, on the formal side, a denotational semantics for RIR
modules (the validation suite currently pins behavior operationally).

\section{Conclusion}
\label{sec:conclusion}

Quantum scientific computing has an algorithm literature that is mature,
growing, and structurally consistent: problems accessed through declared
mathematical input models; solutions assembled from block-encoding algebra
and signal-processing transforms; precision and promises carried as
first-class mathematical objects.  \pqlang{} is a programming language built
to that shape: oracles are declared, carried open, and linked late; solver
families are protocol positions into which interchangeable kernels are
injected; programs remain module-level, symbolic in scale, serializable,
and honest about what they assume and what they omit.  Its coverage---six
ODE solver families, several PDE paths, a reversible-arithmetic and
math-function compiler feeding physics kernels, and a QSVT/QSP engine that
verifies its own synthesis---is, to our knowledge, not matched as an
integrated whole by existing quantum programming systems: several
ecosystems provide parts of the primitive layer (\secref{sec:related}),
but none combines them with open, re-linkable oracle programs and six
ODE solver families.  Its verification discipline is deliberately
graduated, and the campaign of \secref{sec:numverify} shows what that
discipline buys: register-exact certification against independent oracles
for the whole library, two real defects exposed and recorded rather than
silently absorbed, and a clean boundary to what remains open---validation
at scale, the defect repairs and lowering gaps, and the physical layer of
resource estimation, reserved rather than overclaimed.  We offer
the requirements analysis of \secref{sec:requirements} as a reusable lens:
languages for the fault-tolerant scientific-computing era will be judged by
how directly they let the mathematics be written---and by how honestly they
track what remains unproven.
